\documentclass[10pt,a4paper]{amsart}
\usepackage[margin=19mm]{geometry}
\usepackage{amsmath,amssymb,mathtools}
\usepackage[scr=rsfs,cal=euler]{mathalfa}
\usepackage{xfrac}
\usepackage[unicode,hidelinks,bookmarks=false]{hyperref}
\newcommand{\ie}{i.e.\ }
\newcommand{\cf}{cf.\ }
\newcommand{\resp}{respectively\ }
\newcommand{\Subsection}{\subsection}
\providecommand{\nobreakdash}{\nobreak\hbox{-}}
\setlength{\emergencystretch}{2em}
\multlinegap0pt
\usepackage{bm}
\usepackage{bbm}
\usepackage{dsfont}
\usepackage{xspace}
\usepackage{cancel}
\usepackage{mathtools}
\usepackage{amsthm}
\makeatletter
\@ifundefined{theorem}{\newtheorem{theorem}{Theorem}}{}
\@ifundefined{proposition}{\newtheorem{proposition}[theorem]{Proposition}}{}
\makeatother
\DeclareMathOperator{\tr}{tr}
\newcommand{\R}{\mathbb R}
\title[A Lewy theorem for harmonic quasiregular mappings in three-s]{A Lewy theorem for harmonic quasiregular mappings in three-space}
\author{David Kalaj, Jian-Feng Zhu}
\date{Source: arXiv:2607.04720v3 (CC BY 4.0)}
\begin{document}
\maketitle

\noindent\emph{Source and licence.} A Lewy theorem for harmonic quasiregular mappings in three-space. Version arXiv:2607.04720v3 (CC BY 4.0), distributed under the Creative Commons Attribution 4.0 International licence, as recorded in the arXiv licence metadata for this exact version and shown on the abstract page https://arxiv.org/abs/2607.04720v3. Full rights notice: licenses/arxiv-2607.04720-CC-BY-4.0.txt.\par\smallskip

\noindent\textit{Editorial scope.} Counted unit: the Proposition whose source label is \texttt{prop-trace-formula} in the article (source0.tex lines 762--806, source label \texttt{prop-trace-formula}), delivered together with the article's own complete original proof of it (source0.tex lines 808--1181, one \texttt{proof} environment of 374 lines). No mathematics is written, repaired or re-derived: statement and proof are the article's own text line for line. Required context retained uncounted: none. Every definition, display and label of those lines is retained unchanged. Omitted from this package and not redistributed: 1--761, 807--807, 1182--2589. Nothing inside a retained excerpt is omitted or altered: every retained body is delivered exactly as the article's lines are written, so no omission receipt of any kind accompanies this package. Citations: the retained excerpts carry no citation command at all, so no bibliography entry of the article is transcribed and no reference is redistributed. Reference adaptation: none was needed and none was made; every reference inside the delivered unit resolves inside it, so no unresolved reference or citation is produced. Printed slips kept literal: no printed word, symbol, hypothesis or number of the counted statement or of its proof was altered. Layout and class: the article's own class is \texttt{amsart} with packages including inputenc, fontenc, lmodern, geometry, amsmath, amssymb, amsthm, mathtools, graphicx, mathrsfs, enumitem, hyperref, nowidow; the wrapper is the lane's pinned \texttt{amsart} editorial wrapper at 10pt on A4, which restates the theorem-like environments the retained text needs and adds the article's own macro definitions that the retained text needs (\texttt{\textbackslash R}, \texttt{\textbackslash tr}), with extra wrapper packages (bm, bbm, dsfont, xspace, cancel, mathtools). The delivered numbering is the unit's own. Assets: no retained span contains a figure environment, a graphics-inclusion command, a PDF-page inclusion, a table or a TikZ picture; no image file is copied into this package, the package declares no asset dependency and no image file is redistributed with it. Licence: version arXiv:2607.04720v3 of this article is distributed under the Creative Commons Attribution 4.0 International licence, as recorded in the arXiv licence metadata for this exact version and shown on the abstract page https://arxiv.org/abs/2607.04720v3. A full rights notice is delivered at licenses/arxiv-2607.04720-CC-BY-4.0.txt. Characters: every retained excerpt is pure ASCII, so the delivered engine, XeLaTeX with its default Latin Modern text font, reports no missing character.
\begin{proposition}
\label{prop-trace-formula}
Let \(d=n-1\), and let \(P:\R^n\to\R^n\) be homogeneous harmonic of degree
\(m\). Use the local coordinates
\[
        \theta(u)=\bigl(u,\sqrt{1-|u|^2}\bigr)
\]
near \(N=e_n\). Suppose that, after an orientation-preserving linear change
in the target,
\begin{equation}\label{normalization}
        F(0)=e_n,
        \qquad
        F_\alpha(0)=e_\alpha,\qquad \alpha=1,\ldots,d,
\end{equation}
where \(F(u)=P(\theta(u))\). Put
\[
        j(u)=J_P(\theta(u)),
        \qquad
        H^i_{\alpha\beta}=\partial_{\alpha\beta}F^i(0),
\]
\[
        T_\gamma=\sum_{p=1}^{d}H^p_{p\gamma},
        \qquad
        Q(H)=\sum_{\gamma=1}^{d}\sum_{p,q=1}^{d}
              H^p_{q\gamma}H^q_{p\gamma},
        \qquad
        \mu=(m-1)(m+d).
\]
Then
\begin{equation}\label{first-derivative-j}
        \partial_\gamma j(0)=mT_\gamma,
\end{equation}
and
\begin{equation}\label{trace-formula}
        \Delta j(0)
        =
        m\left(
          \sum_{\gamma=1}^{d}T_\gamma^2
          -
          Q(H)
          -
          (d-1)\mu
        \right).
\end{equation}
\end{proposition}

\begin{proof}
We write
\[
        w(u)=\sqrt{1-|u|^2},
        \qquad
        \theta(u)=(u,w(u)).
\]
Thus, \(\theta(0)=e_n\). We also write
\[
        F_\alpha=\partial_\alpha F,
        \qquad
        \theta_\alpha=\partial_\alpha\theta.
\]
Define the \(n\times n\) matrix
\[
        C(u)=\bigl(F_1(u),\ldots,F_d(u),F(u)\bigr),
        \qquad
        s(u)=\det C(u).
\]
The point of this notation is that the Jacobian of \(P\) can be computed by
comparing the frame
\[
        \theta_1,\ldots,\theta_d,\theta
\]
in the domain with its image under \(DP\).

First, we compute the determinant of the domain frame. Since
\[
        \theta_\alpha=(e_\alpha,w_\alpha),
\]
the matrix with columns \(\theta_1,\ldots,\theta_d,\theta\) is
\[
        \begin{pmatrix}
        I_d & u\\
        \nabla w^T & w
        \end{pmatrix}.
\]
Therefore,
\[
        \det(\theta_1,\ldots,\theta_d,\theta)
        =
        w-\nabla w\cdot u.
\]
Because
\[
        \nabla w=-\frac{u}{w},
\]
we get
\[
        w-\nabla w\cdot u
        =
        w+\frac{|u|^2}{w}
        =
        \frac{w^2+|u|^2}{w}
        =
        \frac1w.
\]
Hence,
\[
        \det(\theta_1,\ldots,\theta_d,\theta)=\frac1{w(u)}.
\]
On the other hand,
\[
        DP(\theta(u))\theta_\alpha(u)=F_\alpha(u).
\]
Also, by Euler's identity for a homogeneous map of degree \(m\),
\[
        DP(\theta(u))\theta(u)=mP(\theta(u))=mF(u).
\]
Taking determinants, and using the orientation fixed by the frame above, we
obtain
\[
        J_P(\theta(u))\det(\theta_1,\ldots,\theta_d,\theta)
        =
        \det(F_1,\ldots,F_d,mF).
\]
Thus,
\[
        j(u)\frac1{w(u)}=ms(u),
\]
or equivalently
\begin{equation}\label{j-mws}
        j(u)=mw(u)s(u).
\end{equation}

At \(u=0\), the normalization \eqref{normalization} gives
\[
        F(0)=e_n,\qquad F_\alpha(0)=e_\alpha,\qquad \alpha=1,\ldots,d.
\]
Therefore,
\[
        C(0)=I,
        \qquad
        s(0)=1.
\]
Moreover,
\[
        w(0)=1,
        \qquad
        \nabla w(0)=0,
        \qquad
        \Delta w(0)=-d.
\]
Indeed,
\[
        w(u)=1-\frac12|u|^2+O(|u|^4).
\]
Applying the Laplacian to \eqref{j-mws} and using the product rule, we get
\[
        \frac1m\Delta j(0)
        =
        \Delta(ws)(0)
        =
        w(0)\Delta s(0)+s(0)\Delta w(0)+2\nabla w(0)\cdot\nabla s(0).
\]
Hence,
\begin{equation}\label{Deltaj-Deltas}
        \frac1m\Delta j(0)=\Delta s(0)-d.
\end{equation}

It remains to compute \(\Delta s(0)\). Since \(C(0)=I\), the standard first
and second differential formulas for the determinant at the identity give
\[
        s_\gamma(0)=\tr C_\gamma(0),
\]
and
\begin{equation}\label{det-second}
        s_{\gamma\gamma}(0)
        =
        \tr C_{\gamma\gamma}(0)
        +
        \bigl(\tr C_\gamma(0)\bigr)^2
        -
        \tr\bigl(C_\gamma(0)^2\bigr).
\end{equation}
Indeed,
\[
        D(\det)_I(A)=\tr A,
        \qquad
        D^2(\det)_I(A,A)=(\tr A)^2-\tr(A^2).
\]

We now compute the first trace. Since
\[
        C(u)=\bigl(F_1(u),\ldots,F_d(u),F(u)\bigr),
\]
we have
\[
        C_\gamma(u)=\bigl(F_{1\gamma}(u),\ldots,F_{d\gamma}(u),F_\gamma(u)\bigr).
\]
At \(u=0\), the last column is
\[
        F_\gamma(0)=e_\gamma.
\]
Since \(\gamma\le d\), this last column contributes zero to the last diagonal
entry. Therefore,
\begin{equation}\label{trace-Cgamma}
        \tr C_\gamma(0)
        =
        \sum_{p=1}^{d}F^p_{p\gamma}(0)
        =
        \sum_{p=1}^{d}H^p_{p\gamma}
        =
        T_\gamma.
\end{equation}
Together with \(\nabla w(0)=0\), this gives
\[
        \partial_\gamma j(0)=ms_\gamma(0)=mT_\gamma,
\]
which proves \eqref{first-derivative-j}.

We next compute
\[
        \Delta s(0)=\sum_{\gamma=1}^{d}s_{\gamma\gamma}(0).
\]
From \eqref{det-second},
\[
        \Delta s(0)
        =
        \sum_{\gamma=1}^{d}\tr C_{\gamma\gamma}(0)
        +
        \sum_{\gamma=1}^{d}T_\gamma^2
        -
        \sum_{\gamma=1}^{d}\tr\bigl(C_\gamma(0)^2\bigr).
\]
Thus, we have to evaluate the first and third terms.

Since each coordinate \(F^i\) is the restriction to the sphere of a
homogeneous harmonic polynomial of degree \(m\), it is a spherical harmonic
of degree \(m\). Hence,
\[
        \Delta_{S^d}F^i=-\lambda F^i,
        \qquad
        \lambda=m(m+d-1).
\]
At \(u=0\), the local formula for \(\Delta_{S^d}\) reduces to the Euclidean
tangent Laplacian, and therefore
\begin{equation}\label{second-spherical}
        \sum_{\gamma=1}^{d}F^i_{\gamma\gamma}(0)
        =
        -\lambda F^i(0).
\end{equation}
Differentiating the local expression
\[
        \Delta_{S^d}\phi
        =
        \sum_{\alpha,\beta=1}^{d}
        (\delta_{\alpha\beta}-u_\alpha u_\beta)\phi_{\alpha\beta}
        -
        d\sum_{\alpha=1}^{d}u_\alpha\phi_\alpha
\]
and evaluating at \(u=0\), we obtain
\begin{equation}\label{third-spherical}
        \sum_{\gamma=1}^{d}F^i_{\gamma\gamma\alpha}(0)
        =
        -(\lambda-d)F^i_\alpha(0).
\end{equation}
Indeed, the derivative of the coefficient
\((\delta_{\alpha\beta}-u_\alpha u_\beta)\) vanishes at \(u=0\), while
differentiating the first-order term gives \(-dF^i_\alpha(0)\).

Now
\[
        C_{\gamma\gamma}(u)
        =
        \bigl(F_{1\gamma\gamma}(u),\ldots,F_{d\gamma\gamma}(u),
              F_{\gamma\gamma}(u)\bigr).
\]
Therefore,
\[
        \tr C_{\gamma\gamma}(0)
        =
        \sum_{p=1}^{d}F^p_{p\gamma\gamma}(0)
        +
        F^n_{\gamma\gamma}(0).
\]
Summing over \(\gamma\), and using \eqref{third-spherical} with \(i=p\) and
\(\alpha=p\), we get
\[
        \sum_{\gamma=1}^{d}F^p_{p\gamma\gamma}(0)
        =
        \sum_{\gamma=1}^{d}F^p_{\gamma\gamma p}(0)
        =
        -(\lambda-d)F^p_p(0).
\]
By the normalization \(F_p(0)=e_p\), we have \(F^p_p(0)=1\). Thus,
\[
        \sum_{\gamma=1}^{d}F^p_{p\gamma\gamma}(0)=-(\lambda-d).
\]
Summing this over \(p=1,\ldots,d\), and using \eqref{second-spherical} with
\(i=n\), we obtain
\begin{equation}\label{trace-Cgammagamma}
        \sum_{\gamma=1}^{d}\tr C_{\gamma\gamma}(0)
        =
        -d(\lambda-d)-\lambda.
\end{equation}

It remains to compute the square term. For fixed \(\gamma\), the matrix
\(C_\gamma(0)\) has the block form
\[
        C_\gamma(0)
        =
        \begin{pmatrix}
        A_\gamma & e_\gamma\\
        b_\gamma^T & 0
        \end{pmatrix},
\]
where
\[
        (A_\gamma)^p_q=H^p_{q\gamma},
        \qquad
        (b_\gamma)_q=H^n_{q\gamma},
        \qquad
        p,q=1,\ldots,d.
\]
For such a block matrix,
\[
        \tr
        \begin{pmatrix}
        A & v\\
        b^T & 0
        \end{pmatrix}^{\!2}
        =
        \tr(A^2)+2b^Tv.
\]
Here \(v=e_\gamma\), so
\[
        b_\gamma^Te_\gamma=H^n_{\gamma\gamma}.
\]
Also,
\[
        \tr(A_\gamma^2)=\sum_{p,q=1}^{d}H^p_{q\gamma}H^q_{p\gamma}.
\]
Consequently,
\[
        \tr\bigl(C_\gamma(0)^2\bigr)
        =
        \sum_{p,q=1}^{d}H^p_{q\gamma}H^q_{p\gamma}
        +
        2H^n_{\gamma\gamma}.
\]
Summing over \(\gamma\), we obtain
\[
        \sum_{\gamma=1}^{d}\tr\bigl(C_\gamma(0)^2\bigr)
        =
        Q(H)+2\sum_{\gamma=1}^{d}H^n_{\gamma\gamma}.
\]
But \(H^n_{\gamma\gamma}=F^n_{\gamma\gamma}(0)\), and by
\eqref{second-spherical} with \(i=n\),
\[
        \sum_{\gamma=1}^{d}F^n_{\gamma\gamma}(0)
        =
        -\lambda F^n(0)
        =
        -\lambda,
\]
because \(F(0)=e_n\). Therefore,
\begin{equation}\label{trace-Cgamma-square}
        \sum_{\gamma=1}^{d}\tr\bigl(C_\gamma(0)^2\bigr)
        =
        Q(H)-2\lambda.
\end{equation}

Combining \eqref{det-second}, \eqref{trace-Cgamma},
\eqref{trace-Cgammagamma}, and \eqref{trace-Cgamma-square}, we get
\[
\begin{aligned}
        \Delta s(0)
        &=
        \sum_{\gamma=1}^{d}T_\gamma^2
        +
        \bigl[-d(\lambda-d)-\lambda\bigr]
        -
        \bigl[Q(H)-2\lambda\bigr]  \\
        &=
        \sum_{\gamma=1}^{d}T_\gamma^2
        -
        Q(H)
        -
        (d-1)\lambda
        +
        d^2 .
\end{aligned}
\]
Finally, by \eqref{Deltaj-Deltas},
\[
        \frac1m\Delta j(0)
        =
        \sum_{\gamma=1}^{d}T_\gamma^2
        -
        Q(H)
        -
        (d-1)\lambda
        +
        d^2-d.
\]
Since
\[
        \lambda-d=m(m+d-1)-d=(m-1)(m+d)=\mu,
\]
we obtain
\[
        \Delta j(0)
        =
        m\left(
          \sum_{\gamma=1}^{d}T_\gamma^2
          -
          Q(H)
          -
          (d-1)\mu
        \right),
\]
which is \eqref{trace-formula}.
\end{proof}

\end{document}
